% !TeX root = ../thesis.tex

\chapter{结构设计与强度计算}

\label{chap:struct-design}

本章在第二章确定的设计基础和材料参数之上，完成承压壳体的结构设计与强度验证。结构设计和强度计算本为一体两面——结构构型的选型决定了计算模型的边界条件，而强度准则的计算结果又反过来约束结构方案的可行性。本章将筒体、封头、裙座等主要承压零部件的结构设计与其强度验证一体化呈现，以工程实践中"设计-校核-迭代"的真实逻辑链展开阐述。

\section{筒体设计}

\subsection{筒体结构方案}

冷媒膨胀罐的筒体采用单层卷焊式圆筒形结构，由一张或两张 Q345R 钢板经冷卷或温卷成形后焊接纵缝，再与上、下封头组焊环缝。该方案在 DN1800 直径级别上工艺成熟，制造成本低，焊缝质量控制手段完备\cite{nbt47013-2015}。

筒体基本参数确定如下：内径取工艺数据表给定值 $D_i = 1800$ mm；切线长度（T/T，不含封头直边段）$L = 1700$ mm，使罐体长径比 $L/D_i \approx 0.94$，属于短粗型立式容器，有利于气液分离和自由液面的稳定；名义厚度 $\delta_n = 14$ mm 的确定过程见本节强度计算。

\subsection{内压强度计算}

筒体计算厚度按照 GB/T 150.3—2024 式(5-1)确定。该式基于薄壁圆筒中径公式，假定环向薄膜应力沿壁厚均匀分布，适用于 $D_o/D_i \le 1.5$（本设备 $D_o/D_i = 1828/1800 = 1.016$，远小于准则限值）。

\begin{equation}
\delta = \frac{P_c D_i}{2[\sigma]^t \phi - P_c}
\label{eq:cyl-delta}
\end{equation}

式中：$P_c$——计算压力，$0.80$ MPa（液柱静压力小于设计压力的 $5\%$，按 GB/T 150.1—2024 第 3.1.7 条可忽略）；$D_i$——筒体内径，$1800$ mm；$[\sigma]^t$——设计温度许用应力，$186$ MPa；$\phi$——焊接接头系数，$0.85$。

代入计算：$\delta = 0.80 \times 1800 / (2 \times 186 \times 0.85 - 0.80) = 1440 / 315.4 = 4.57$ mm。

计算厚度 $4.57$ mm 仅为内压强度准则的基本需求。在此基础上，设计厚度（增加腐蚀裕量）$\delta_d = 4.57 + 3.0 = 7.57$ mm，进一步考虑钢板负偏差后名义厚度至少需 $\delta_n \ge 7.87$ mm。但内压需求所对应的壁厚远不能覆盖全真空外压失稳的刚度需求，且难以满足开孔补强对壳体多余厚度的依赖和制造运输工况的截面刚度要求。综合外压稳定性（见第 \ref{sec:ext-pressure} 节）、开孔补强（见第 \ref{sec:opening} 节）和结构刚度的多重约束，取名义厚度 $\delta_n = 14$ mm。

有效厚度 $\delta_e = \delta_n - C_1 - C_2 = 14 - 0.3 - 3.0 = 10.7$ mm。该值将用于后续所有应力校核计算。

操作工况下筒体环向薄膜应力的校核结果：$\sigma_\theta = P_c(D_i+\delta_e)/(2\delta_e) = 67.7$ MPa，远低于许用限值 $[\sigma]^t\phi = 158.1$ MPa，安全裕度 $n_\theta = 2.33$。轴向薄膜应力 $\sigma_z = P_c(D_i+\delta_e)/(4\delta_e) = 33.9$ MPa，安全裕度 $n_z = 4.66$。内压操作工况下筒体的强度储备充裕，壁厚由刚度条件主导。

\section{外压稳定性}

\label{sec:ext-pressure-2}

\subsection{失稳机制与计算方法}

设计条件要求容器具备全真空运行能力。在全真空工况下，圆筒承受 $P_{\mathrm{vac}} = 0.1013$ MPa 的均匀外压。外压圆筒的失效模式并非强度屈服，而是弹性或弹塑性失稳——当外压力达到临界值时，圆筒截面突然丧失原始圆形而发生压扁或褶皱。GB/T 150.3—2024 第 6 章提供了基于小变形弹性稳定性理论的图算法，适用于本设备 $D_o/\delta_e = 170.8$ 位于弹性失稳区间的工况。

\subsection{几何参数与许用外压}

外压计算长度取筒体切线长度与两端封头曲面当量直段之和。封头对外压筒体起刚性约束作用，椭圆封头曲面的当量直段长度按 $h_i/3$ 计入：

\begin{align}
L_{cr} &= L + 2 \times \frac{h_i}{3} = 1700 + 2 \times \frac{450}{3} = 2000 \; \mathrm{mm} \\
D_o    &= D_i + 2\delta_n = 1828 \; \mathrm{mm} \\
\frac{D_o}{\delta_e} &= \frac{1828}{10.7} = 170.8 \\
\frac{L_{cr}}{D_o} &= \frac{2000}{1828} = 1.094
\end{align}

由 $L_{cr}/D_o = 1.094$ 和 $D_o/\delta_e = 170.8$ 查取 GB/T 150.3—2024 图 6-2，得外压应变系数 $A = 1.6 \times 10^{-4}$。$A$ 值位于材料弹性模量控制段，许用外压力直接由弹性模量 $E_t$ 确定：

\begin{equation}
[P] = \frac{2AE_t}{3(D_o/\delta_e)}
\label{eq:allow-p-2}
\end{equation}

全真空工况金属温度 $100$ ℃，Q345R 弹性模量保守取 $E_t \approx 2.0 \times 10^5$ MPa。代入得：$[P] = 2 \times 1.6 \times 10^{-4} \times 2.0 \times 10^5 / (3 \times 170.8) = 0.125$ MPa。

\subsection{稳定性评定}

$[P] = 0.125 \; \mathrm{MPa} > P_{\mathrm{vac}} = 0.1013 \; \mathrm{MPa}$，筒体在全真空工况下不会发生外压失稳。外压稳定安全系数 $n_s = [P]/P_{\mathrm{vac}} = 1.23 > 1.0$，满足 GB/T 150.3—2024 要求。

\section{封头设计}

上、下封头均采用标准椭圆形封头（EHA 型，按 GB/T 25198—2010 选型\cite{gbt25198-2010}），以内径为基准，$D_i/2h_i = 2$，形状系数 $K = 1.0$。封头曲面深度 $h_i = D_i/4 = 450$ mm，直边段高度 $h_s = 40$ mm（标准值），单片封头总高 $H_f = h_i + h_s = 490$ mm。

封头计算厚度按 GB/T 150.3—2024 式(7-1)确定：

\begin{equation}
\delta_h = \frac{K P_c D_i}{2[\sigma]^t \phi - 0.5P_c}
\label{eq:head-delta-2}
\end{equation}

代入计算得 $\delta_h = 4.56$ mm。封头计算厚度与筒体计算厚度 $4.57$ mm 几乎一致——这不是巧合，而是标准椭圆封头与等壁厚圆筒内压承载的力学匹配所决定的。封头的设计厚度 $\delta_{dh} = 4.56 + 3.0 = 7.56$ mm。

封头名义厚度取与筒体等厚 $\delta_{nh} = 14$ mm，主要考虑：成型工艺减薄（冲压工艺约 $10\%\sim15\%$ 壁厚损失，需投料厚度 $\ge 16$ mm）；与筒体等厚有利于采购和环缝焊接；封头曲面上开设多个接管开孔后对壳体强度的削弱。封头成型后实测最小厚度不得低于 $\delta_{eh} = 10.7$ mm，此为设计有效厚度的硬性下限。

封头最大薄膜应力校核：在封头顶部（曲率半径最大区域），薄膜应力 $\sigma_h \approx 67.5$ MPa $ < [\sigma]^t\phi = 158.1$ MPa，满足强度要求。

\section{裙座设计}

\subsection{选型与结构}

立式容器采用圆筒形焊制裙座支承，裙座设计标准为 NB/T 47065.5—2018\cite{nbt47065-2018}，强度校核参照 NB/T 47041—2014\cite{nbt47041-2014}。裙座筒节外径 $D_{so} = 1828$ mm，与筒体外径相等，壁厚初选 $\delta_{ss} = 14$ mm（与壳体同厚，简化采购），筒节长度 $H_s \approx 2000$ mm。裙座上部与下封头直边段外侧以全焊透对接环缝连接，下部焊接基础环板。

裙座腐蚀裕量取 $C_{2s} = 2.0$ mm（大气腐蚀 + 涂层老化环境），小于壳体腐蚀裕量。裙座不与工艺介质接触，受力以轴向压缩（自重 + 风/地震弯矩）和横向剪切为主。裙座壳体应力、基础环板弯曲应力及地脚螺栓拉应力的详细校核按 NB/T 47041—2014 执行，属本设计范围。

\subsection{附件与开孔}

裙座侧壁开设：上部通气孔 $\ge \phi 80$ mm（或等面积长圆孔），对称布置，数量不少于 2 个，防止有害气体积聚；底部排液孔 $\ge \phi 20$ mm，不少于 2 个；检修人孔 $\ge \phi 450$ mm，便于人员进入检查基础螺栓。基础环板外径大于裙座外径，厚度 $\ge 20$ mm，地脚螺栓数量和规格由风/地震工况下的倾覆力矩校核确定（本设计阶段暂不展开）。

裙座外表面涂装：环氧富锌底漆（$\ge 60$ μm）+ 环氧云铁中间漆（$\ge 100$ μm）+ 聚氨酯面漆（$\ge 60$ μm），按 HG/T 20584—2020\cite{hgt20584-2020} 及 HG/T 20580—2020\cite{hgt20580-2020} 附录要求执行。该涂层体系在连云港沿海大气环境中的预期有效防护年限为 $10\sim15$ 年。

\section{接管与法兰}

\subsection{接管选型与布置}

工艺接管与仪表接管的规格、位置和功能汇总于表 \ref{tab:nozzle-specs}。接管材质为 20 号钢无缝钢管（GB 6479—2013）\cite{gb6479-2013}，管壁厚度由内压圆筒公式计算后圆整为标准壁厚系列，接管计算压力与壳体相同，许用应力按 20 号钢管材在 170 ℃ 下的取值（$\approx 140$ MPa）执行。

\begin{table}[!htb]
  \centering
  \caption{接管参数一览}
  \label{tab:nozzle-specs}
  \small\renewcommand{\arraystretch}{1.25}
  \setlength{\tabcolsep}{5pt}
  \begin{tabularx}{\textwidth}{@{}c>{\centering\arraybackslash}X>{\centering\arraybackslash}X>{\centering\arraybackslash}Xc>{\centering\arraybackslash}X@{}}
    \toprule
    \textbf{编号} & \textbf{用途} & \textbf{所在部件} & \textbf{规格 $\phi \times \delta_t$} & \textbf{DN} & \textbf{外伸/内伸 $L_o$/$L_i$ (mm)} \\
    \midrule
    N1   & 冷媒进口    & 上封头顶   & $\phi 114.3 \times 6.0$ & 100 & 200 / 0 \\
    N2   & 氮气补气口  & 上封头     & $\phi 60.3 \times 5.0$  & 50  & 200 / 0 \\
    N3   & 放空口      & 上封头     & $\phi 60.3 \times 5.0$  & 50  & 200 / 0 \\
    N4   & 压力表口    & 上封头     & $\phi 33.7 \times 4.0$  & 25  & 150 / 0 \\
    N5   & 温度计口    & 上封头     & $\phi 48.3 \times 5.0$  & 40  & 250 / 400 \\
    N6a  & 液位计口(上)& 筒体侧壁  & $\phi 33.7 \times 4.0$  & 25  & 200 / 0 \\
    N6b  & 液位计口(下)& 筒体侧壁  & $\phi 33.7 \times 4.0$  & 25  & 200 / 0 \\
    N7   & 压力变送器口 & 筒体侧壁  & $\phi 33.7 \times 4.0$  & 25  & 150 / 0 \\
    N8   & 冷媒出口    & 下封头底   & $\phi 114.3 \times 6.0$ & 100 & 250 / 0 \\
    N9   & 排净口      & 下封头底   & $\phi 60.3 \times 5.0$  & 50  & 200 / 0 \\
    M1   & 人孔        & 筒体侧壁   & $\phi 508 \times 12$     & 500 & 300 / 0 \\
    \bottomrule
  \end{tabularx}
\end{table}

注：N5 温度计口内伸 $400$ mm，使测温元件插入至罐体半径约 1/2 深度处，确保温度测量值代表主体介质温度而非壁面附近边界层温度。

\subsection{法兰选型}

所有接管法兰按 HG/T 20592～20635—2009 选型\cite{hgt20592-2009}，选用带颈对焊法兰（WN），密封面型式为突面（RF）。选型依据：WN 法兰的颈部锥形过渡可降低法兰环与接管连接处的应力集中，在内压—全真空交变工况和温度波动条件下抗疲劳性能优于板式平焊法兰；法兰与接管为对接焊连接，焊缝可实施射线或超声检测，质量可控。

法兰紧固件选用等长双头螺柱配六角螺母，材质 35CrMoA（调质），垫片为柔性石墨复合垫或 PTFE 包覆垫，按 HG/T 20606—2009\cite{hgt20606-2009} 和 HG/T 20613—2009\cite{hgt20613-2009} 选取。法兰压力等级的选择原则：设计压力 $0.80$ MPa 远低于 PN16（1.6 MPa）的公称压力基准，各法兰均按 PN16 等级选取可有充裕承压裕量。压力表口和温度计口因仪表制造商标准接口要求，选用 PN25 等级。

\subsection{人孔}

在筒体侧壁开设 DN500 回转盖带颈对焊法兰人孔（M1），按 HG/T 21514—2014 选取\cite{hgt21514-2014}。人孔参数：公称直径 DN500，公称压力 PN16，法兰密封面 RF，接管尺寸 $\phi 508 \times 12$，人孔盖为回转盖式。人孔开孔直径为全设备最大开孔（$d = 484$ mm），其补强设计是本章开孔补强校核的核心对象。

\subsection{防涡流挡板}

罐底冷媒出口（N8）处设置十字形防涡流挡板，材质 Q345R，焊接固定于出口接管上方的筒体内壁。挡板的功能为破坏排液时液面旋涡，防止气体进入出流管道导致气蚀或流量失稳。挡板形式和尺寸参考 HG/T 20580—2020 附录的推荐做法\cite{hgt20580-2020}。

\section{开孔补强}

\label{sec:opening}

\subsection{补强方法与准则}

筒体及封头上的接管开孔破坏了壳体的连续承载截面，孔边产生应力集中，峰值应力可达壳体薄膜应力的 $2\sim3$ 倍。按 GB/T 150.3—2024 第 8 章\cite{gbt150-2024}，采用等面积补强法（Area Replacement Method），其判据为：在有效补强范围内，补强金属的截面积 $A_e$ 不小于因开孔而削弱的壳体金属截面积 $A$。

补强有效范围：宽度方向取 $B = \max\{2d, d + 2\delta_n + 2\delta_{nt}\}$，沿壳体外壁的外侧高度 $h_1 = \min\{\sqrt{d\delta_{nt}}, \text{接管实际外伸}\}$，内侧高度 $h_2 = \min\{\sqrt{d\delta_{nt}}, \text{接管实际内伸}\}$。补强面积为以下各部分之和：壳体在补强范围内的多余厚度面积 $A_1$、接管在补强范围内的多余厚度面积 $A_2$、以及焊缝面积 $A_3$。

\subsection{补强计算结果}

各主要开孔的补强计算结果汇总于表 \ref{tab:reinforcement}。计算结果表明：DN500 和 DN100 开孔区域的接管厚壁自补强已足以满足等面积补强准则，无需增设补强圈；DN50 及以下小口径仪表接管因开孔直径小，壳体有效厚度提供的多余金属已覆盖削弱面积，不需另行补强。

\begin{table}[!htb]
  \centering
  \caption{开孔补强计算结果}
  \label{tab:reinforcement}
  \small\renewcommand{\arraystretch}{1.2}
  \setlength{\tabcolsep}{6pt}
  \begin{tabularx}{\textwidth}{@{}c>{\centering\arraybackslash}X>{\centering\arraybackslash}X>{\centering\arraybackslash}X>{\centering\arraybackslash}X>{\centering\arraybackslash}X@{}}
    \toprule
    \textbf{序号} & \textbf{开孔位置} & \textbf{接管规格} & \textbf{$A$ / mm²} & \textbf{$A_e$ / mm²} & \textbf{补强方式} \\
    \midrule
    1 & 筒体   & $\phi 508 \times 12$ & 2212 & 3647 & 厚壁接管自补强 \\
    2 & 筒体   & $\phi 114.3 \times 6$ & 468 & 1020 & 厚壁接管自补强 \\
    3 & 上封头 & $\phi 114.3 \times 6$ & 468 & 980  & 厚壁接管自补强 \\
    4 & 下封头 & $\phi 114.3 \times 6$ & 468 & 1020 & 厚壁接管自补强 \\
    5 & 上/下封头 & $\phi 60.3 \times 5$ & 230 & 465 & 不需另行补强 \\
    6 & 多位置   & $\phi 33.7 \times 4$ & 117 & 380 & 不需另行补强 \\
    \bottomrule
  \end{tabularx}
\end{table}

\section{设备总尺寸链}

设备总体尺寸通过以下几何关系链确定，是 CAD 总装图绘制的基准数据：

壳体总高（不含裙座）：$H_{\mathrm{shell}} = L + 2H_f = 1700 + 2 \times 490 = 2680$ mm。裙座高度（暂定）$H_s \approx 2000$ mm。设备总高（近似）$H_{\mathrm{total}} = H_{\mathrm{shell}} + H_s \approx 4680$ mm。

罐体满液液柱高度：$H_{\mathrm{liquid}} = L + 2h_i = 1700 + 2 \times 450 = 2600$ mm，对应满液介质质量约 $5500$ kg。基础顶面（TOS）标高 $\approx 9050$ mm。

上述尺寸均已在 CAD 总装图中完整表达，并作为各零部件图（筒体图、封头图、裙座图、管口方位图）的基准定位依据。
